% Generated by tools/build_v6.py; edit the assembly rules, not this file.
\section{One-unit null-work obstructions}\label{sec:lower-foundations}

This appendix supplies the sign-testing and annular ingredients used in
Section~\ref{sec:matching-lower}. The weighted sign-testing mechanism is
an adaptation of Chen--Li--Qiao's gap-entropy argument, including their
Bernoulli extension \cite[Lemma 4.1 and Remark D.4]{clq}. We give the full
finite-gap version needed here, followed by the scale-sensitive gate and
its actual minimum-band-mass constraint.

\subsection{Weighted sign testing}

\begin{lemma}[Weighted sign entropy]\label{lem:weighted-sign}
Let $d_1>\cdots>d_K>0$, with $d_j/d_{j+1}\ge2$ and $d_1\le1/128$.
One observable binary sign test, with the same independent seed law at all
gaps, has error at most $1/10$ under every Bernoulli law $1/2\pm d_j$.
Suppose its stopping time $\tau$ has finite expectation under both signs at
every such gap.  Put $c_j=d_j^2\E_{-d_j}\tau$.
For every probability vector $\pi$ on these gaps,
\begin{equation}\label{eq:sc-weighted-sign}
 \sum_{j=1}^K\pi_jc_j\ge
 \frac{1+\operatorname{Ent}(\pi)}{214}.
\end{equation}
The stopping and decision rules are the same measurable functions of the
seed and observation tape at every gap; the realized gap is not supplied
as an input. The test may depend on the entire known finite gap list and
on the weights. Zero weights are allowed, with $0\log(1/0)=0$.
No termination or correctness assumption under the fair coin is required.
\end{lemma}

\begin{proof}
First, the ordinary stopped change-of-measure inequality gives a uniform
positive floor for each $c_j$.  The terminal sign decision has minus/plus
KL divergence at least
\[
 \operatorname{kl}(9/10,1/10)=\frac45\log9>\frac85.
\]
One Bernoulli observation has divergence at most
\[
 \frac{(2d_j)^2}{(1/2+d_j)(1/2-d_j)}\le\frac{64}{3}d_j^2.
\]
The stopped chain rule, justified by finite expected endpoint work, implies
$c_j>3/40>1/16$. More explicitly, apply the chain rule first to
the transcript capped at $N$. Encode a nonterminal decision by a third
symbol. Its probability tends to zero at both endpoints, so lower
semicontinuity of finite-alphabet KL and monotone convergence of
$\E_{-d_j}(\tau\wedge N)$ justify the stopped inequality.
This step uses no fair-law stopping assumption.

Set $b_0=1/1024$ and
\[
 n_j=\left\lfloor b_0/d_j^2\right\rfloor,\qquad
 u_j=\left\lceil4\E_{-d_j}\tau\right\rceil.
\]
Since $d_j\le1/128$, we have $n_j\ge b_0/(2d_j^2)\ge8$.
The KL divergence between the complete prefixes of length $n_j$ at the
minus and plus laws is at most $1/48$.  These prefixes include the
independent seed and unused bits if the test has already stopped.
Pinsker's inequality therefore gives total variation less than $1/8$.
An early minus is a plus-law error, so its probability under the minus law
is at most $1/10+1/8$.  Markov's inequality gives
$\Pp_{-d_j}(\tau>u_j)\le1/4$.  Consequently the event
\[
 E_j=\{\text{output minus},\ n_j<\tau\le u_j\}
\]
has $\Pp_{-d_j}(E_j)>1/4$.

Let $\Pp_0$ denote the fair-coin law with the same seed and policy.
Comparing deterministic prefixes of length $u_j$ gives KL at most
$4d_j^2u_j\le16c_j+1$.  Binary data processing and
$\operatorname{kl}(p,z)\ge p\log(1/z)-\log2$ yield
\begin{equation}\label{eq:sc-window-mass}
 \Pp_0(E_j)\ge\exp(-64c_j-8).
\end{equation}
Only deterministic finite prefixes have been compared with the fair law.

We next bound a weighted overlap pathwise.  On a path with finite terminal
time $t$, participation in $E_j$ implies
\[
 t>n_j\ge\frac{b_0}{2d_j^2},\qquad
 t\le4\E_{-d_j}\tau+1.
\]
In particular $t\ge2$ and
$c_j\ge(t-1)d_j^2/4\ge td_j^2/8$.
Order the participating gaps from smallest to largest.  Their squared
values grow by at least a factor of four, and the smallest has
$td_j^2>b_0/2$.  The corresponding costs are thus bounded below by
$b_0 4^r/16=4^{r-7}$, for $r=0,1,\ldots$.  Hence
\[
 \sum_j\frac{\1_{E_j}}{1+c_j}
 \le\sum_{r\ge0}\frac1{1+4^{r-7}}
 \le7+\sum_{r\ge7}4^{7-r}=\frac{25}{3}.
\]
On a nonterminating fair path no $E_j$ occurs, so the same inequality
holds.  Integrating and using \eqref{eq:sc-window-mass}, together with
$1+c_j\le e^{c_j}$, gives
\[
 Z:=\sum_je^{-65c_j}\le(25/3)e^8,
 \qquad \log Z<11.
\]
Gibbs' inequality for $\pi$ and the normalized masses $e^{-65c_j}/Z$
implies
\[
 \operatorname{Ent}(\pi)\le65\sum_j\pi_jc_j+\log Z.
\]
Writing $\bar c=\sum_j\pi_jc_j$ and retaining the stronger floor
$\bar c\ge3/40$, we obtain
\[
 1+\operatorname{Ent}(\pi)
 \le65\bar c+9+\log(25/3)
 \le\left[65+\frac{40}{3}\{9+\log(25/3)\}\right]\bar c
 <214\bar c.
\]
The windows need not be disjoint; only the displayed weighted-overlap
bound is used.
\end{proof}


\subsection{A moment comparison on the half-prior class}

Write $X_\rho=\{0\le q\le1/2:\int q\,d\mu\le\rho\}$.
For an integrable observable one-unit output $Y$ with finite uniform
variance and uniform bias at most $\rho/64$, put $V=\sup_{q\in X_\rho}\Var_qY$ and
$M=\sup_{q\in X_\rho}\E_qY^2$.

\begin{lemma}[Moment comparison]\label{lem:lf-moment}
For $0<\rho\le1/4$, one has $V\ge2\rho/3$ and $M\le3V/2$.
\end{lemma}
\begin{proof}
Let $f(d)=\E_{d,+}Y-\E_{d,-}Y$. Comparing the null population with
the constant prior $q=\rho$ gives $\int f\,d\mu\ge31/32$.
Conditional variance at this constant prior and Jensen's inequality yield
\[
 V\ge\rho(1-\rho)\int f^2\,d\mu
 \ge\rho(1-\rho)(31/32)^2\ge2\rho/3.
\]
Every mean has absolute value at most $65\rho/64$, so
$M\le V+(65\rho/64)^2\le3V/2$.
The constant prior also gives integrability of both fixed-sign second
moments; all conditional statements may be taken outside a common null set.
\end{proof}

\subsection{An activation gate that respects the magnitude scale}

The gate argument charges activation at the information
scale \(a^2\E_0T\).  In this subsection \(\mu\) is any probability law
supported on \((0,a]\), \(a\le1/4\), and
\[
 \mathcal X_\rho=
 \{q:0\le q\le\tfrac12,\ \textstyle\int q\,d\mu\le\rho\}.
\]
Let \((Y,T)\) be an observable one-coin procedure with uniform bias at
most \(\rho/64\), finite uniform variance \(V\), and finite expected work
under every \(q\in\mathcal X_\rho\).  Write
\[
 M_0=\sup_{q\in\mathcal X_\rho}\E_qY^2,\qquad
 W_0=\E_0T,\qquad P=M_0W_0.
\]
Here $M_0$ is the uniform second-moment bound, whereas $W_0$ is
null work; the subscript on $M_0$ does not restrict its supremum to the null.
The moment comparison in Lemma~\ref{lem:lf-moment} gives
\(M_0\le3V/2\).

More explicitly, observability means that an independent seed $U$ and
the genuine observations generate $\mathcal F_n=\sigma(U,X_1,\ldots,X_n)$,
that $T$ is a nonnegative integer-valued stopping time, and that $Y$ is
$\mathcal F_T$-measurable. The procedure has no access to the latent
magnitude or sign. Restarts below refresh $U$ and the unused observation
block, while retaining that same latent state. These are the observation-model assumptions of
Section~\ref{sec:rare-model}.

\begin{lemma}[A scale-sensitive baseline and gated activation]
\label{lem:information-gate}
Put
\[
 C_g=e^7,\qquad \kappa=64/3,\qquad
 c_g=(256\kappa C_g^4)^{-1}.
\]
Every procedure above satisfies
\begin{equation}
 P\ge\frac{\rho}{64a^2}.
 \label{eq:cross-baseline}
\end{equation}
If \(a^2P\le c_g\), there is an observable procedure \(Z\) on the same
one-coin experiment, with moment bound \(m_Z\), null work \(W_0^Z\), and
a measurable set \(\mathcal G\) of good magnitudes, such that
\begin{equation}
 \begin{gathered}
 \sup_{q\in\mathcal X_\rho}\E_qZ^2\le m_Z,\qquad
 \rho/16\le m_Z\le2C_g^4a^2P,\\
 m_ZW_0^Z\le C_g^4P,\qquad
 \mu(\mathcal G^c)\le144m_Z.
 \end{gathered}
 \label{eq:cross-gate-bounds}
\end{equation}
At every good gap, both signs have finite conditional expected work and
\begin{equation}
 \left|\E[Z\mid d,0]\right|\le\frac18,\qquad
 \E[Z\mid d,1]\ge\frac12,\qquad
 \E[Z^2\mid d,0]\le\frac{m_Z}{\rho},\qquad
 \E[Z^2\mid d,1]\le\frac{3m_Z}{\rho}.
 \label{eq:cross-gate-good}
\end{equation}
One common sign test with both errors below \(1/10\) on these gaps uses
\begin{equation}
 n=\left\lceil1024m_Z/\rho\right\rceil
 \le1040m_Z/\rho
 \label{eq:cross-gate-restarts}
\end{equation}
independent restarts of \(Z\) on unused observations of the same coin.
Restart independence is conditional on its fixed latent magnitude and
sign, not after averaging over the population.
Its negative expected work is \(n\,\E[T_Z\mid d,0]\).
\end{lemma}

\begin{proof}
We first prove \eqref{eq:cross-baseline}.  Let \(P_-\) and \(P_+\) be
the full transcript laws with uniformly negative and uniformly positive
signs, respectively.  The priors \(q=0\) and \(q\equiv\rho\) give
\[
 \E_+Y-\E_-Y\ge31/32,\qquad
 \E_-Y^2\le M_0,\qquad \E_+Y^2\le M_0/\rho.
\]
The positive moment and expected work are finite because \(q\equiv\rho\)
is admissible.  Revealing \(D\) for the information calculation, and
using the stopped chain rule, gives
\[
 \mathfrak h(P_-,P_+)\le\tfrac12\kappa a^2W_0,
\]
where \(\mathfrak h=1-\operatorname{Aff}\).
Indeed the one-observation endpoint divergence is at most
\(\kappa d^2\) for \(d\le1/4\).
For any square-integrable output \(F\),
\begin{equation}
 |\E_PF-\E_QF|^2
 \le4(\E_PF^2+\E_QF^2)\mathfrak h(P,Q).
 \label{eq:cross-hellinger-moment}
\end{equation}
Thus
\[
 (31/32)^2
 \le(5\kappa/2)a^2P/\rho,
\]
which implies \eqref{eq:cross-baseline}.

The decision to take a first observation is seed-only and precedes
all data. The information gate below is instead an observable
finite-prefix stopping rule; neither step reveals a latent magnitude.
For the gate, on a full observation tape define
\[
 S_n=\sum_{i=1}^n(2X_i-1),\qquad
 n_0=\lceil a^{-2}\rceil,\qquad H_0=\lceil a^{-1}\rceil,
\]
\begin{equation}
 G=\min\{n_0,\inf\{n\ge1:|S_n|\ge H_0\}\},\qquad
 \sigma=T\wedge G,\qquad E=\{T\ge G\}.
 \label{eq:cross-gate}
\end{equation}
The event \(E\) is measurable at \(\sigma\), including the tie \(T=G\).
Include the entire independent seed in the stopped prefix.
At \(\sigma\), \(|S_\sigma|\le H_0\) and \(\sigma\le n_0\).
For a fixed \(d\le a\) and sign \(s\), the likelihood relative to the
fair stopped-prefix law is
\[
 \exp\left\{(2s-1)S_\sigma\operatorname{arctanh}(2d)
       +\frac{\sigma}{2}\log(1-4d^2)\right\}.
\]
The inequalities
\[
 \operatorname{arctanh}(2d)\le8a/3,\qquad
 -\tfrac12\log(1-4d^2)\le8a^2/3
\]
show that the absolute logarithm is at most \(37/6<7\).
Every fixed-state prefix law, and the all-negative magnitude mixture,
therefore has density in \([C_g^{-1},C_g]\) relative to fair.

Let \(p=\Pp_{\rm fair}(E)\) and \(p_s(d)=\Pp_{d,s}(E)\).
The number $p$ is determined by the known policy, its seed law, and a
bounded fair-prefix experiment. It does not depend on the unknown prior
or genuine latent state. This is an existence argument; no efficient
computation of arbitrary policy probabilities is claimed. No termination
under the fair law beyond this bounded prefix is needed. Then
\begin{equation}
 p/C_g\le p_s(d)\le C_gp.
 \label{eq:cross-gate-probabilities}
\end{equation}
Bounded fair stopping gives
\(\E_{\rm fair}S_\sigma^2=\E_{\rm fair}\sigma\).
A gate activation either reaches \(n_0\) or hits \(H_0\), so
\begin{equation}
 p\le\left(n_0^{-1}+H_0^{-2}\right)\E_{\rm fair}\sigma
 \le2a^2\E_{\rm fair}\sigma
 \le2C_ga^2W_0.
 \label{eq:cross-gate-cost}
\end{equation}

Set \(Y_e=Y\1_{\{T<G\}}\), explicitly zero on \(E\).
This is a function of the bounded prefix.
Its second moment under either fixed state is at most \(C_g^2M_0\).
The stopped endpoint divergence and the same prefix domination give
\[
 \KL(P_{-d}^{\sigma},P_{+d}^{\sigma})
 \le\kappa d^2\E_{-d}\sigma
 \le\kappa C_g^2a^2W_0.
\]
Since \(\mathfrak h\le\KL/2\), \eqref{eq:cross-hellinger-moment} implies
\begin{equation}
 |\E_{+d}Y_e-\E_{-d}Y_e|^2
 \le4\kappa C_g^4a^2P.
 \label{eq:cross-early-output}
\end{equation}
Under \(a^2P\le c_g\), every such early mean difference is at most
\(1/8\) in absolute value.  In this branch \(p>0\): if \(p=0\), the
gate never activates under any state, so \(Y=Y_e\), contradicting
the integrated mean separation \(31/32\) from \(q=0,q\equiv\rho\).

Here is an actual observable procedure realizing the conditioning.
Keep the same genuine coin.  Repeatedly run the original procedure,
using a fresh original seed and unused iid flips for each trial.
Reject a trial if its original stopping time is before its gate.
When a trial reaches the gate, accept it, continue that original trial
to its original stopping time, and return
\[
 Z=C_gpY.
\]
Every rejected observation is charged.  Conditional on \(d,s\), trials
are iid with acceptance probability \(p_s(d)>0\).
Writing \(w_s(d)=\E[T\mid d,s]\), the renewal identity gives
\begin{equation}
 \E[T_Z\mid d,s]=\frac{w_s(d)}{p_s(d)}.
 \label{eq:cross-renewal}
\end{equation}
This identity includes the work of all rejected trials and the final
accepted trial. To verify it without assuming independence of acceptance
and trial cost, let $L$ be the first accepted trial. Conditional on $(d,s)$,
\[
 \E\sum_{k=1}^{L}T_k
 =\sum_{k\ge1}\E[T_k\mathbf1_{\{L\ge k\}}]
 =\sum_{k\ge1}(1-p_s(d))^{k-1}w_s(d)
 =w_s(d)/p_s(d).
\]
The event $\{L\ge k\}$ involves only preceding trials, so it is independent
of the current trial cost. It is not necessary that $T_k$ be independent
of its own acceptance event. Moreover,
\[
 \E[Z^2\mid d,s]
 =\frac{C_g^2p^2}{p_s(d)}\E[Y^2\1_E\mid d,s]
 \le C_g^3p\,\E[Y^2\mid d,s].
\]
With \(m_Z=C_g^3pM_0\), integration and
\eqref{eq:cross-gate-cost}--\eqref{eq:cross-renewal} give
\begin{equation}
 \sup_q\E_qZ^2\le m_Z,\qquad
 m_Z\le2C_g^4a^2P,\qquad
 W_0^Z\le C_gW_0/p,\qquad
 m_ZW_0^Z\le C_g^4P.
 \label{eq:cross-gate-moments}
\end{equation}
Expected work under every \(q\) is finite by the same comparison.

We establish conditional sign separation directly; conditioning has
not preserved the estimator's means.  Write
\[
 f_s(d)=\E[Y\mid d,s],\quad
 \ell_s(d)=\E[Y\1_E\mid d,s],\quad
 z_s(d)=\E[Z\mid d,s].
\]
The set \(\{f_+-f_-<3/4\}\) has mass below \(2\rho\).
Otherwise, writing that set as $S$ with $\mu(S)\ge2\rho$, the
prior $q=(\rho/\mu(S))\mathbf1_S$ lies in $\mathcal X_\rho$ and has total
mass $\rho$. It would contradict the bias guarantee and the mean
separation $31\rho/32$. This uses fractional prior values, so it applies
unchanged when $\mu$ has atoms.
Outside that set, \eqref{eq:cross-early-output} gives
\(\ell_+-\ell_-\ge5/8\).
Also
\[
 z_s(d)=\frac{C_gp}{p_s(d)}\ell_s(d),\qquad
 1\le\frac{C_gp}{p_s(d)}\le C_g^2.
\]
If \(|z_-(d)|\le1/8\), then \(|\ell_-(d)|\le1/8\),
\(\ell_+(d)\ge1/2\), and \(z_+(d)\ge1/2\).
The exceptional mass for \(|z_-|>1/8\) is at most \(64m_Z\).
Since
\[
 m_Z\le2C_g^4c_g=1/(128\kappa)<1/256,
\]
this preliminary mean-good set has mass at least
\(1-2\rho-64m_Z\ge1/4\).
Taking \(q\equiv\rho\) in the second-moment bound yields
\(m_Z\ge\rho/16\).

Finally, put \(a_Z(d)=\E[Z^2\mid d,0]\) and
\(b_Z(d)=\E[Z^2\mid d,1]\).
The set \(\{a_Z>m_Z/\rho\}\) has mass at most \(\rho\).
The set \(\{b_Z-a_Z>2m_Z/\rho\}\) has mass below \(2\rho\),
again by using a fractional admissible prior on a set of mass
at least \(2\rho\).  Removing these sets and null sets of infinite
conditional work gives total exceptional mass
\[
 5\rho+64m_Z\le144m_Z
\]
and proves \eqref{eq:cross-gate-good}.
Average \(n\) independent restarts of \(Z\) on unused blocks of the
same coin, and threshold at \(5/16\).
Both mean margins are at least \(3/16\), so Chebyshev gives errors
at most
\[
 \frac{3m_Z/\rho}{n(3/16)^2}\le1/12<1/10.
\]
The bound on \(n\) follows from \(m_Z\ge\rho/16\), and work adds,
proving \eqref{eq:cross-gate-restarts}.
\end{proof}

\subsection{A scaled annular obstruction}

For a positive vector \(v=(v_i)\), set
\[
 \mathcal G(v)=A[1+\Ent(v/A)],\qquad A=\sum_i v_i.
\]
A useful elementary identity is
\begin{equation}
 \frac{\partial\mathcal G}{\partial v_i}
 =1+\log(A/v_i)\ge1.
 \label{eq:cross-entropy-monotone}
\end{equation}
Thus this cost-entropy functional is coordinatewise increasing.
In particular, if \(u\le v_i\le U\) for \(K\) coordinates, then
\begin{equation}
 uK(1+\log K)\le\mathcal G(v)\le UK(1+\log K).
 \label{eq:cross-entropy-sandwich}
\end{equation}

\begin{proposition}[A gated annular lower bound]
\label{prop:scaled-annular}
Suppose \(\mu\) is supported on \((0,a]\), \(a\le1/4\).
Take a nonempty finite annular family
\[
 I_j=(b_j/2,b_j],\qquad b_j/b_{j+1}\ge4,\qquad
 b_1\le1/128,
\]
with positive masses \(w_j=\mu(I_j)\).  Define
\[
 w_*=\min_jw_j,\qquad
 A=\sum_jw_j/b_j^2,\qquad
 \pi_j=\frac{w_j/b_j^2}{A},\qquad
 \Gamma^{\circ}=A[1+\Ent(\pi)].
\]
Every one-coin procedure with the assumptions above satisfies
\begin{equation}
 V\E_0T\ge
 \frac{e^{-28}}{2\,000\,000}
       \min\{\rho\Gamma^{\circ},w_*/a^2\}.
 \label{eq:cross-scaled-annular}
\end{equation}
\end{proposition}

\begin{proof}
If \(a^2P\ge c_g\), use \(w_*\le1\) and \(M_0\le3V/2\).
Otherwise apply Lemma~\ref{lem:information-gate}.
If \(m_Z>w_*/288\), then
\[
 P>w_*/(576C_g^4a^2).
\]
If \(m_Z\le w_*/288\), at least half the mass of every annulus is
good.  Averaging selects \(d_j\in I_j\cap\mathcal G\) with
\[
 w_-^Z(d_j)\le\frac4{w_j}\int_{I_j}w_-^Z(d)\,\mu(dd).
\]
The selected gaps are separated by a factor greater than two and
bounded by \(1/128\). The representative points may depend on the
procedure being bounded. This is legitimate: the weighted sign lemma
holds for every such finite list. After the points have been selected,
the same already-constructed restart test is run at all of them; it is
not retuned to the selected point. The weights remain the normalized
band costs $\pi_j$, not the original band probabilities.
Apply Lemma~\ref{lem:weighted-sign} to the common restart test,
with weights \(\pi_j\):
\[
 \frac{1+\Ent(\pi)}{214}
 \le n\sum_j\pi_jd_j^2w_-^Z(d_j)
 \le\frac{4n}{A}W_0^Z.
\]
Consequently,
\[
 \rho\Gamma^{\circ}
 \le856\rho nW_0^Z
 \le890\,240\,m_ZW_0^Z
 \le890\,240\,C_g^4P.
\]
Converting \(P=M_0W_0\) to \(VW_0\) completes all branches with the
displayed conservative constant.  No atomlessness is used.
\end{proof}

The entropy branch is available when $\rho a^2\Gamma^{\circ}\le w_*$.
This condition is invariant under a common rescaling of all magnitudes.
